% Author: Prof. Dr. Matthias Jung, DL9MJ
% Year: 2025
\begin{tikzpicture}
    \begin{axis}[%
        /pgf/number format/1000 sep={ },
        /pgf/number format/use comma,
        axis lines=middle,
        axis line style={-Triangle},
        ytick={-2,-1,0,1,2},
        every major tick/.append style={thick, black},
        yticklabels={-2~V,-1~V,0,1~V,2~V},
        grid=major,
        grid style={line width=.1pt, draw=gray!10},
        major grid style={line width=.2pt,draw=gray!50},
        xticklabel=\empty,
        xtick style={draw=none},
        xmin=0,
        xmax=64,
        ymin=-2.4,
        ymax=3.0,
        xmajorgrids=false,
        xlabel style={at={(ticklabel* cs:1)},anchor=north east},
        ylabel style={at={(ticklabel* cs:1)},anchor=north east},
        ylabel={$U_\mathrm{A}$},
        xlabel={$t$},
        width=\linewidth*\getDarcImageFactor,
        height=4cm,
        scale only axis,
    ]

    % konstante Taktzeitpunkte
    \addplot[
        only marks,
        mark=*,
        mark size=0.45pt,
        gray,
        samples at={0,1,...,64}
    ] {-2.2};

    % Phaseninkrement +1
    % N = 32 Tabelleneinträge, pro Takt wird die Phase um 1 erhöht
    \addplot[
        const plot,
        samples at={0,1,...,64},
        DARCblue,
        thick
    ] {
        (round(sin(360*x/32) * 1.5/2 * 7.5 + 0.5) - 0.5) / 7.5 * 2
    };

    \node[anchor=north west, DARCblue] at (axis cs:2,2.7) {$\varphi_{n+1} = \varphi_n + 1$};

    \end{axis}
\end{tikzpicture}%